\section{Study Assessments and Procedures}

\draftinstr{Follow the Phase 2 publication section
\texttt{final-protocol/publication/sections/sec-08-assessments.tex} and the Phase 1
model file \texttt{trial-protocol/draft-protocol/sections/sec-08-assessments.tex}.
State that every centrally readable endpoint is read centrally across the eight
sites (central imaging core BICR, central masked pathology, central laboratories).
Timing is governed by the SoA (\S1.3). Source from
\texttt{trial-protocol/nih-protocol/06\_study\_assessments\_and\_procedures.md}.}

\subsection{Efficacy Assessments}
\draftinstr{Define PFS (RECIST 1.1, investigator-assessed primary with BICR
sensitivity, \S9.4.2) and OS to 24 months \cite{DutchCohort2025}; the central
masked R0 resection rate and MPR; ISGPS biochemical-leak / Grade B / Grade C
fistula grading \cite{Bassi2017}; week-12 KRAS ctDNA clearance \cite{Tie2019ctdna};
the Arm A anastomosis-quality reads against the ring-tension bands; and QoL via
EORTC QLQ-C30 and PAN26 \cite{EORTCqlqc30}.}

\subsection{Safety and Other Assessments}
\draftinstr{Clinical safety (vitals, exam, hematology, chemistry, coagulation, the
Arm A serum daraxonrasib trough, imaging) plus the Arm A continuous device-telemetry
stream under 21 CFR \S312.23(g) (per-arm and cumulative force profiles, trajectory
accuracy, vascular-zone verdicts, heartbeat / watchdog events, e-stop activations).
Render \texttt{trial-phase-2/mermaid/fig-15-vascular-safety-gate.md} as a TikZ
\texttt{mermaidfig} (this is \textbf{Figure 15}): the five-vessel vascular
safety-zone gate (soft-warning 3.0 mm, no-fly 1.0/1.5 mm, hard-stop 5.0 mm) sampled
at 10 kHz. Then render \texttt{trial-phase-2/mermaid/fig-16-estop-architecture.md}
as a TikZ \texttt{mermaidfig} (this is \textbf{Figure 16}): the heartbeat /
watchdog / e-stop architecture ($100\,\mu$s watchdog, $50\,\mu$s arm park, $\leq 3$
ms cross-arm e-stop, $\leq 500$ ms hardware e-stop). Source from
\texttt{trial-protocol/inputs/2030-pdac-1min-final-paper/sections/methods.tex} and
\texttt{trial-protocol/inputs/21cfr312\_adapt/05\_clinical\_holds\_appendices\_closing.tex}.}

\subsection{Adverse Events and Serious Adverse Events}
\draftinstr{Two parallel binding streams: conventional pharmacovigilance for both
arms and a Physical AI safety stream for Arm A. Render
\texttt{trial-phase-2/mermaid/fig-17-ae-reporting.md} as a TikZ \texttt{mermaidfig}
(this is \textbf{Figure 17}): the combined AE / Physical AI AE reporting workflow
(CTCAE v5 grading, 7-day fatal / 15-day serious-unexpected timelines, the six
Physical AI triggers, audit-trail preservation $-24$ h to $+72$ h federated across
sites, and notification of FDA / IRB / DSMB / Safety Review Committee). Source from
\texttt{trial-protocol/inputs/21cfr312\_adapt/03\_protocol\_amendments\_reporting.tex}.}

\subsubsection{Definition of Adverse Events (AE)}
\draftinstr{Standard AE definition, CTCAE v5 grading in both arms, plus the
Physical AI adverse event per 21 CFR \S312.32(f) (robotic malfunction, AI
prediction error, sensor failure, communication loss, unauthorized access,
digital-twin divergence, any e-stop activation) \cite{kawchak2026adapt312}.}

\subsubsection{Definition of Serious Adverse Events (SAE)}
\draftinstr{Standard SAE criteria plus the serious Physical AI adverse event per
\S312.32(f) (potential SAE outcome, unplanned interruption of drug administration,
uncontrolled drug release, or sterile-field compromise).}

\subsubsection{Classification of an Adverse Event}
\draftinstr{Four axes: severity (CTCAE v5); separate drug and device causality
(procedure causality in Arm B); expectedness against the reference safety
information; and seriousness. Unexpected + serious + at least possibly related
defines a reportable suspected adverse reaction.}

\subsubsection{Time Period and Frequency for Assessment and Follow-Up}
\draftinstr{Solicit AEs from consent; intensive window through the 30-day visit;
SAEs, events of special interest, Physical AI AEs, and new anticancer therapy
collected to the 24-month OS endpoint.}

\subsubsection{Adverse Event Reporting}
\draftinstr{Record all AEs in the CRF with onset / resolution, grade, drug and
device causality, expectedness, seriousness, action, and outcome; Arm A telemetry-
derived observations cross-referenced to audit-trail records.}

\subsubsection{Serious Adverse Event Reporting}
\draftinstr{Investigator-to-sponsor within 24 hours; FDA timelines (7-day fatal /
life-threatening, 15-day other serious unexpected) plus the six Physical AI
reporting triggers under \S312.32(g); preserve the hash-chained audit trail $-24$ h
to $+72$ h, federated across the eight sites \cite{ecfr2026part11}.}

\subsubsection{Events of Special Interest}
\draftinstr{Vascular injury (including any vascular-breach hard-stop e-stop),
clinically relevant fistula (ISGPS B/C \cite{Bassi2017}), conversion, any
cybersecurity incident, and any model-performance degradation; reviewed by the
Physical AI Safety Review Committee ($\leq 90$-day cadence) and reported to the
DSMB.}

\subsubsection{Reporting of Pregnancy}
\draftinstr{Pregnancy reported (not an AE) within 24 hours and followed to outcome;
highly effective contraception required; an associated SAE reported on the expedited
timeline.}

\subsection{Unanticipated Problems}

\subsubsection{Definition}
\draftinstr{Standard unanticipated-problem definition; treat the Physical AI hold
grounds (repeated uncorrected safety failures, confirmed cybersecurity compromise,
simulation-validation deficiency) as paradigm unanticipated problems.}

\subsubsection{Reporting}
\draftinstr{Report to the single IRB and, as applicable, the sponsor and FDA on the
7- or 15-day SAE timeline or promptly otherwise; device-related reports preserve the
federated audit trail \cite{ecfr2026part11}; the Physical AI Safety Review Committee
and DSMB review each.}

\subsubsection{Reporting to Participants}
\draftinstr{Communicate to affected participants via IRB-approved re-consent or
written notification; update the consent and re-consent enrolled participants when
the risk profile changes, consistent with the Physical AI opt-out.}
